\chapter{Referencia del protocolo IPC}
\label{anx:ipc}

Este anexo documenta el contrato técnico completo de la comunicación entre el orquestador Java y los motores Python, ampliando la descripción del protocolo incluida en la sección~\ref{sec:ipc}.

\section*{Estructura del mensaje}

Todo intercambio sigue la estructura de \textit{framing} binario definida en \texttt{IpcMessagePackCodec} y \texttt{ipc\_protocol.py}: un prefijo de cuatro bytes en orden de red que indica la longitud del cuerpo, seguido del cuerpo serializado en MessagePack. El cuerpo corresponde siempre a un objeto envolvente con los cuatro campos definidos en el registro \texttt{IpcEnvelope}:

\begin{table}[h]
\centering
\caption{Campos del objeto envolvente \texttt{IpcEnvelope}.}
\label{tab:ipc-envelope}
\begin{tabular}{@{}lll@{}}
\toprule
\textbf{Campo} & \textbf{Tipo} & \textbf{Descripción} \\
\midrule
\texttt{protocol\_version} & \texttt{String} & Versión del protocolo; actualmente \texttt{"1.0"} \\
\texttt{message\_type}     & \texttt{String} & Código semántico de la operación (véase tabla~\ref{tab:ipc-tipos}) \\
\texttt{correlation\_id}   & \texttt{String} & Identificador UUID único por petición \\
\texttt{payload}           & objeto          & Contenido específico de la operación \\
\bottomrule
\end{tabular}
\end{table}

\section*{Tipos de mensaje}

La enumeración \texttt{IpcMessageType} define los veintiún códigos de operación reconocidos por los motores actuales. La tabla~\ref{tab:ipc-tipos} los agrupa por dominio funcional.

\begin{table}[h]
\centering
\caption{Tipos de mensaje del protocolo IPC, agrupados por dominio.}
\label{tab:ipc-tipos}
\begin{tabular}{@{}lll@{}}
\toprule
\textbf{Dominio} & \textbf{Tipo} & \textbf{Dirección} \\
\midrule
Descarga de datos
    & \texttt{FETCH\_REQUEST}      & Java $\rightarrow$ Python \\
    & \texttt{FETCH\_CHUNK}        & Python $\rightarrow$ Java (respuesta parcial) \\
    & \texttt{FETCH\_RESPONSE}     & Python $\rightarrow$ Java (fin de descarga) \\
\midrule
Indicadores técnicos
    & \texttt{INDICATORS\_REQUEST}  & Java $\rightarrow$ Python \\
    & \texttt{INDICATORS\_RESPONSE} & Python $\rightarrow$ Java \\
\midrule
Entrenamiento de modelos
    & \texttt{TRAIN\_REQUEST}       & Java $\rightarrow$ Python \\
    & \texttt{TRAIN\_RESPONSE}      & Python $\rightarrow$ Java \\
\midrule
Optimización de hiperparámetros
    & \texttt{OPTIMIZE\_REQUEST}    & Java $\rightarrow$ Python \\
    & \texttt{OPTIMIZE\_RESPONSE}   & Python $\rightarrow$ Java \\
\midrule
Inferencia
    & \texttt{PREDICT\_REQUEST}     & Java $\rightarrow$ Python \\
    & \texttt{PREDICT\_RESPONSE}    & Python $\rightarrow$ Java \\
\midrule
Simulación histórica
    & \texttt{BACKTEST\_REQUEST}    & Java $\rightarrow$ Python \\
    & \texttt{BACKTEST\_RESPONSE}   & Python $\rightarrow$ Java \\
\midrule
Trading en tiempo real
    & \texttt{RT\_SIGNAL}           & Python $\rightarrow$ Java (señal de orden) \\
    & \texttt{RT\_LOG}              & Python $\rightarrow$ Java (línea de registro) \\
\midrule
Control
    & \texttt{ERROR}                & Python $\rightarrow$ Java \\
\bottomrule
\end{tabular}
\end{table}

\section*{Convención de códigos de salida}

Complementando los tipos de mensaje, los motores Python comunican el resultado final de su ejecución mediante el código de retorno del proceso, que el orquestador Java interpreta para decidir si la operación fue exitosa o requiere tratamiento de error:

\begin{table}[h]
\centering
\caption{Códigos de salida de los motores Python.}
\label{tab:exit-codes}
\begin{tabular}{@{}cl@{}}
\toprule
\textbf{Código} & \textbf{Significado} \\
\midrule
0 & Ejecución completada con éxito \\
1 & Error de lógica no recuperable sin intervención del usuario \\
  & (estrategia inválida, modelo no entrenado, etc.) \\
2 & Error en los datos de entrada \\
3 & Interrupción por tiempo de espera agotado o cierre externo \\
\bottomrule
\end{tabular}
\end{table}

\section*{Comportamiento de lectura y escritura}

La función de lectura del módulo \texttt{ipc\_protocol.py} implementa tres niveles de compatibilidad para garantizar interoperabilidad durante transiciones de versión: en primer lugar intenta leer el mensaje como un envolvente con prefijo de longitud de cuatro bytes (protocolo v1); si la longitud declarada no concuerda con el tamaño del flujo recibido, reintenta la deserialización directamente como MessagePack sin prefijo (modo legado); y como último recurso interpreta el contenido como texto JSON plano. La función de escritura construye siempre el envolvente completo en formato v1, añadiendo un identificador de correlación único generado en tiempo de ejecución, de modo que todas las respuestas de todos los motores siguen el mismo formato con independencia de la operación solicitada.
